\section{Job Object}

The Job object represents the root of a MIFX package.

It provides the entry point for interpreting the machining process and defines the primary structure of the package.

The Job object is stored in the file:

\begin{verbatim}
job.json
\end{verbatim}

located at the root of the MIFX archive.

\subsection{Job Structure}

A Job object typically contains:

\begin{itemize}
\item package metadata
\item setup references
\item tool group references
\item optional addon references
\end{itemize}

Example:

\begin{verbatim}
{
  "mifxVersion": "0.1",

  "header": {
    "export": {
      "id": "d347ed61-2e99-4b69-8946-fcba14292330",
      "at": "2026-03-11T19:03:24Z",
      "by": "Fusion360",
      "exporter": "MIFXExporter"
    },
    "job": {
      "name": "3+2-Milling v5"
    }
  },

  "setups": [
    {
      "id": "set-1",
      "path": "setups/set-1/setup.json"
    }
  ],

  "tools": [
    {
      "role": "tool_group",
      "kind": "json",
      "path": "tools/main.json",
      "present": true
    }
  ]
}
\end{verbatim}

\subsection{MIFX Version}

The \texttt{mifxVersion} field specifies the version of the MIFX specification used to create the package.

Example:

\begin{verbatim}
"mifxVersion": "0.1"
\end{verbatim}

Software implementations should verify compatibility before processing the package.

\subsection{Header}

The \texttt{header} object provides metadata describing the origin of the package.

Typical information may include:

\begin{itemize}
\item exporting system
\item export timestamp
\item source design document
\item job name
\end{itemize}

The structure of the header is intentionally flexible and may vary between implementations.

\subsection{Setup References}

The \texttt{setups} array lists the setups contained in the package.

Each entry contains:

\begin{itemize}
\item the setup identifier
\item the relative path to the setup object
\end{itemize}

Example:

\begin{verbatim}
{
  "id": "set-1",
  "path": "setups/set-1/setup.json"
}
\end{verbatim}

The order of setups in this array represents the logical order of setups in the machining process.

\subsection{Tool Groups}

Tool definitions are organized into tool groups located in the \texttt{tools} directory.

The Job object references tool groups through artifacts.

Example:

\begin{verbatim}
{
  "role": "tool_group",
  "kind": "json",
  "path": "tools/main.json",
  "present": true
}
\end{verbatim}

Tool groups allow multiple machine configurations (such as milling spindles or lathe turrets) to be represented within the same job.

\subsection{Optional Addons}

The Job object may reference optional addon definitions that extend the package with additional functionality.

Addons are typically stored in the \texttt{addons} directory.

Implementations should ignore addon definitions they do not recognize.